\documentclass[aspectratio=169,11pt]{beamer}

\usetheme{AnnArbor}
\usepackage{fontspec}
\usepackage{amsmath,amssymb,booktabs,array,multicol}
\usepackage{listings}
\usepackage{xcolor}
\usepackage{graphicx}
\usepackage{hyperref}
\usepackage{ragged2e}

\setmainfont{DejaVu Serif}
\setsansfont{DejaVu Sans}
\setmonofont{Noto Sans Mono}

\definecolor{CodeBg}{RGB}{248,248,248}
\definecolor{TitleBlue}{RGB}{38,79,129}
\definecolor{AccentGreen}{RGB}{35,125,93}
\definecolor{SoftGray}{RGB}{80,80,80}
\definecolor{QuizOrange}{RGB}{180,95,30}

\setbeamercolor{title}{fg=TitleBlue}
\setbeamercolor{frametitle}{fg=TitleBlue}
\setbeamercolor{structure}{fg=AccentGreen}
\setbeamercolor{block title}{fg=white,bg=TitleBlue}
\setbeamercolor{block body}{bg=blue!3}
\setbeamercolor{alertblock title}{fg=white,bg=QuizOrange}
\setbeamercolor{alertblock body}{bg=orange!6}
\setbeamertemplate{navigation symbols}{}
\setbeamertemplate{footline}[frame number]
\setbeamertemplate{itemize item}{\raise0.15ex\hbox{\tiny$\blacktriangleright$}}
\setbeamertemplate{itemize subitem}{\raise0.15ex\hbox{\tiny$\bullet$}}

\lstdefinestyle{py}{
  language=Python,
  basicstyle=\ttfamily\scriptsize,
  backgroundcolor=\color{CodeBg},
  frame=single,
  framerule=0.2pt,
  rulecolor=\color{gray!35},
  breaklines=true,
  showstringspaces=false,
  keywordstyle=\color{TitleBlue}\bfseries,
  commentstyle=\color{SoftGray}\itshape,
  stringstyle=\color{AccentGreen},
  upquote=true
}

\newcommand{\key}[1]{\textbf{\textcolor{AccentGreen}{#1}}}
\newcommand{\smallnote}[1]{\vspace{0.25em}{\footnotesize\color{SoftGray}#1}}
\newcommand{\quiztitle}[1]{\textbf{\textcolor{QuizOrange}{#1}}}
\newcommand{\cmd}[1]{\texttt{\detokenize{#1}}}

\title[Giới thiệu Pandas]{Giới thiệu Pandas}
\subtitle{DataFrame, Series, làm sạch dữ liệu, gom nhóm và phân tích dữ liệu kinh doanh}
\author{Duc-Minh Vu}
\institute{SLSCM Lab -- FDA -- NEU}
\date{\today}

\begin{document}

\begin{frame}
  \titlepage
\end{frame}

\begin{frame}{Mục tiêu học tập}
\small
Sau khi hoàn thành bài học, người học có thể:
\begin{itemize}
  \item Giải thích vai trò của Pandas trong khoa học dữ liệu và phân tích kinh doanh.
  \item Phân biệt \cmd{Series} và \cmd{DataFrame}.
  \item Tạo, kiểm tra, chọn, lọc và sắp xếp dữ liệu dạng bảng.
  \item Đọc và ghi dữ liệu CSV, Excel, JSON và tệp văn bản.
  \item Phát hiện và xử lý giá trị thiếu, dòng trùng lặp và kiểu dữ liệu không nhất quán.
  \item Biến đổi cột và áp dụng các hàm đơn giản.
  \item Sử dụng \cmd{groupby()}, \cmd{agg()}, \cmd{merge()}, \cmd{concat()} và pivot table.
  \item Tính thống kê mô tả, tương quan và một số tổng hợp chuỗi thời gian đơn giản.
\end{itemize}
\end{frame}


\begin{frame}{Lộ trình bài học}
\begin{columns}[T]
\column{0.48\linewidth}
\begin{enumerate}
  \item Vì sao cần học Pandas?
  \item Pandas là gì?
  \item Pandas trong quy trình Data Science
  \item Thiết lập môi trường
  \item DataFrame đầu tiên
  \item Series và DataFrame
  \item Đọc lệnh Pandas
\end{enumerate}
\column{0.48\linewidth}
\begin{enumerate}
  \setcounter{enumi}{7}
  \item Kiểm tra, chọn, lọc và sắp xếp dữ liệu
  \item Đọc và ghi tệp dữ liệu
  \item Làm sạch dữ liệu
  \item Tạo cột, gom nhóm và tổng hợp
  \item Merge, concat và pivot table
  \item Tương quan, trực quan hóa, chuỗi thời gian
  \item Bài tập thực hành
\end{enumerate}
\end{columns}
\end{frame}


\section{Pandas là gì?}


\begin{frame}{Vì sao cần học Pandas?}
\small
Pandas là cầu nối thực hành giữa dữ liệu thô và các mô hình khoa học dữ liệu.
\begin{itemize}
  \item Dữ liệu thực tế thường lộn xộn, thiếu giá trị và nằm ở nhiều định dạng khác nhau.
  \item Pandas giúp tổ chức bản ghi thành dòng và thuộc tính thành cột.
  \item Pandas hỗ trợ kiểm tra, làm sạch, biến đổi, tổng hợp và báo cáo dữ liệu.
  \item Pandas làm việc tốt với NumPy, Matplotlib, Scikit-learn và Statsmodels.
\end{itemize}
\begin{block}{Quy trình phân tích kinh doanh}
Dữ liệu thô $\rightarrow$ bảng dữ liệu sạch $\rightarrow$ insight tổng hợp $\rightarrow$ dữ liệu sẵn sàng cho mô hình.
\end{block}
\end{frame}

\begin{frame}{Pandas là gì?}
\small
\key{Pandas} là một thư viện Python mã nguồn mở dùng cho \key{xử lý và phân tích dữ liệu}.
Thư viện này được xây dựng trên NumPy và cung cấp các công cụ cấp cao để làm việc với dữ liệu có cấu trúc và dữ liệu dạng bảng.
\vspace{0.4em}

Pandas đặc biệt hữu ích cho:
\begin{itemize}
  \item đọc dữ liệu từ CSV, Excel, JSON và tệp văn bản;
  \item làm sạch, biến đổi và chuẩn bị dữ liệu;
  \item chọn, lọc, gom nhóm và tổng hợp các quan sát;
  \item kết hợp nhiều tập dữ liệu;
  \item tính thống kê mô tả và trực quan hóa nhanh.
\end{itemize}
\end{frame}


\begin{frame}{Pandas trong quy trình Khoa học dữ liệu}
\small
Pandas thường nằm ở bước chuyển dữ liệu thô thành dữ liệu sạch và có thể phân tích được.
\vspace{0.3em}

\begin{center}
\begin{tabular}{c c c}
\fbox{CSV / Excel / Database} & $\rightarrow$ & \fbox{Pandas DataFrame} \\[0.4em]
 & $\searrow$ & $\downarrow$ \\[0.2em]
\fbox{Raw data} & $\rightarrow$ & \fbox{Clean table} \\[0.4em]
 & $\searrow$ & $\downarrow$ \\[0.2em]
\fbox{Business question} & $\rightarrow$ & \fbox{Insight / Model} \\
\end{tabular}
\end{center}

\begin{itemize}
  \item NumPy hỗ trợ tính toán số học hiệu quả.
  \item Pandas tổ chức dữ liệu theo bảng: dòng, cột và chỉ số.
  \item Matplotlib/Seaborn trực quan hóa kết quả từ DataFrame.
  \item Scikit-learn thường nhận dữ liệu đã được xử lý từ Pandas.
\end{itemize}
\begin{block}{Ý tưởng chính}
Pandas là cầu nối giữa tệp dữ liệu thực tế và phân tích dữ liệu bằng Python.
\end{block}
\end{frame}


\begin{frame}[fragile]{Yêu cầu nền tảng và thiết lập môi trường}
\begin{block}{Kiến thức nên có}
\begin{itemize}
  \item Python cơ bản: biến, list, dictionary, vòng lặp và hàm.
  \item Kiến thức NumPy cơ bản là một lợi thế.
  \item Môi trường làm việc: Jupyter Notebook, JupyterLab, Google Colab, VS Code hoặc công cụ tương tự.
\end{itemize}
\end{block}
\begin{block}{Cài đặt và import}
\begin{lstlisting}[style=py]
pip install pandas
\end{lstlisting}
\begin{lstlisting}[style=py]
import pandas as pd
print(pd.__version__)
\end{lstlisting}
\end{block}
\end{frame}

\begin{frame}[fragile]{Tạo một DataFrame}
\small
\begin{lstlisting}[style=py]
import pandas as pd

data = {
    "Name": ["An", "Binh", "Chi"],
    "Age": [20, 21, 19],
    "GPA": [3.2, 3.6, 3.4]
}

df = pd.DataFrame(data)
print(df)
\end{lstlisting}
\begin{block}{Diễn giải}
Mỗi khóa trong dictionary trở thành một cột. Mỗi list cung cấp các giá trị cho cột đó.
\end{block}
\end{frame}

\begin{frame}[fragile]{Kết quả chạy code: DataFrame đầu tiên}
\small
Khi in \cmd{df}, Pandas hiển thị bảng gồm 3 dòng và 3 cột.
\begin{lstlisting}[style=py]
   Name  Age  GPA
0    An   20  3.2
1  Binh   21  3.6
2   Chi   19  3.4
\end{lstlisting}
\begin{block}{Điểm cần chú ý}
Cột bên trái \cmd{0, 1, 2} là index mặc định do Pandas tự tạo.
\end{block}
\end{frame}

\begin{frame}{Cấu trúc của DataFrame}
\small
\begin{tabular}{p{0.24\linewidth}p{0.64\linewidth}}
\toprule
\textbf{Thành phần} & \textbf{Ý nghĩa} \\
\midrule
Dòng & Bản ghi, quan sát, sinh viên, khách hàng, sản phẩm, đơn hàng. \\
Cột & Biến, thuộc tính, trường dữ liệu, đặc trưng. \\
Index & Nhãn dòng dùng để định danh và chọn dữ liệu. \\
Nhãn cột & Tên dùng để truy cập các biến. \\
Giá trị & Dữ liệu thực tế nằm bên trong bảng. \\
\bottomrule
\end{tabular}
\begin{block}{Ví dụ}
Một bảng sinh viên có thể gồm các cột \cmd{StudentID}, \cmd{Name}, \cmd{Age}, \cmd{Major} và \cmd{GPA}.
\end{block}
\end{frame}


\begin{frame}[fragile]{Series và DataFrame từ ví dụ đầu tiên}
\small
\begin{columns}[T]
\column{0.49\linewidth}
\textbf{Một cột thường là Series}
\begin{lstlisting}[style=py]
print(sales["Price"])
\end{lstlisting}
\begin{lstlisting}[style=py]
0    10
1    20
2     8
Name: Price, dtype: int64
\end{lstlisting}
\column{0.49\linewidth}
\textbf{Nhiều cột là DataFrame}
\begin{lstlisting}[style=py]
print(sales[["Product", "Price"]])
\end{lstlisting}
\begin{lstlisting}[style=py]
  Product  Price
0       A     10
1       B     20
2       C      8
\end{lstlisting}
\end{columns}
\begin{alertblock}{Ghi nhớ}
Một cột có nhãn chỉ số là \cmd{Series}; một bảng hai chiều là \cmd{DataFrame}.
\end{alertblock}
\end{frame}


\begin{frame}{Cách đọc lệnh Pandas}
\small
Phần lớn lệnh Pandas có một trong hai dạng sau:
\[
\cmd{pd.function(input, option=value)} \quad \text{hoặc} \quad \cmd{df.method(option=value)}
\]
\begin{itemize}
  \item \cmd{pd}: bí danh chuẩn của thư viện Pandas.
  \item \cmd{df}: biến thường dùng để lưu một DataFrame.
  \item \cmd{function}: hàm cấp thư viện, ví dụ \cmd{pd.read_csv()}.
  \item \cmd{method}: thao tác gắn với một DataFrame hoặc Series.
  \item \cmd{option=value}: tham số đặt tên dùng để điều khiển hành vi của lệnh.
\end{itemize}
\begin{block}{Ví dụ}
\cmd{pd.read_csv("sales.csv")} đọc tệp CSV và đưa dữ liệu vào một DataFrame.
\end{block}
\end{frame}

\begin{frame}{Mẫu lệnh: đối tượng, thuộc tính, phương thức, hàm}
\small
\begin{tabular}{p{0.27\linewidth}p{0.62\linewidth}}
\toprule
\textbf{Dạng lệnh} & \textbf{Ý nghĩa} \\
\midrule
\cmd{pd.DataFrame(data)} & Gọi hàm Pandas để tạo một bảng dữ liệu. \\
\cmd{df.shape} & Đọc thuộc tính mô tả số dòng và số cột. \\
\cmd{df.head()} & Gọi phương thức gắn với DataFrame hiện có. \\
\cmd{pd.merge(left,right)} & Gọi hàm Pandas để kết hợp hai bảng dữ liệu. \\
\cmd{df["Sales"].mean()} & Chọn một cột rồi gọi phương thức của Series. \\
\bottomrule
\end{tabular}
\vspace{0.5em}
\begin{alertblock}{Ghi chú giảng dạy}
Với mỗi lệnh, sinh viên nên xác định ba thành phần: đầu vào, thao tác và đầu ra.
\end{alertblock}
\end{frame}

\begin{frame}[fragile]{Ví dụ nhỏ: giải thích từng dòng}
\small
\begin{lstlisting}[style=py]
import pandas as pd

sales = pd.DataFrame({
    "Product": ["A", "B", "C"],
    "Quantity": [2, 3, 5],
    "Price": [10, 20, 8]
})

sales["Revenue"] = sales["Quantity"] * sales["Price"]
print(sales)
\end{lstlisting}
\begin{itemize}
  \item \cmd{pd.DataFrame(...)} tạo một tập dữ liệu dạng bảng.
  \item Mỗi khóa trong dictionary trở thành tên cột.
  \item \cmd{sales["Revenue"] = ...} tạo một cột tính toán mới.
  \item Phép nhân được thực hiện theo từng dòng trên hai cột đã chọn.
\end{itemize}
\end{frame}

\begin{frame}[fragile]{Kết quả chạy code: ví dụ doanh thu}
\small
Sau khi chạy ví dụ ở slide trước, cột \cmd{Revenue} được tạo bằng \cmd{Quantity * Price}.
\begin{lstlisting}[style=py]
  Product  Quantity  Price  Revenue
0       A         2     10       20
1       B         3     20       60
2       C         5      8       40
\end{lstlisting}
\begin{block}{Cách đọc kết quả}
Mỗi dòng là một sản phẩm. Ví dụ, sản phẩm B có \cmd{Quantity = 3}, \cmd{Price = 20}, nên \cmd{Revenue = 60}.
\end{block}
\end{frame}

\begin{frame}{Quick Check: Pandas là gì?}
\small
\quiztitle{Câu 1.} Bí danh chuẩn thường dùng cho Pandas là gì?
\begin{multicols}{2}
\begin{itemize}
  \item A. \cmd{pn}
  \item B. \cmd{pd}
  \item C. \cmd{ps}
  \item D. \cmd{pa}
\end{itemize}
\end{multicols}
\vspace{0.6em}
\quiztitle{Câu 2.} Cấu trúc dữ liệu nào của Pandas là hai chiều?
\begin{multicols}{2}
\begin{itemize}
  \item A. \cmd{Series}
  \item B. \cmd{tuple}
  \item C. \cmd{DataFrame}
  \item D. \cmd{ndarray}
\end{itemize}
\end{multicols}
\end{frame}

\section{Cơ bản về DataFrame}


\begin{frame}[fragile]{Kiểm tra nhanh DataFrame}
\small
\begin{lstlisting}[style=py]
print(df.head())
print(df.tail())
print(df.shape)
print(df.columns)
print(df.index)
print(df.dtypes)
df.info()
print(df.describe())
\end{lstlisting}
\begin{block}{Mục đích}
Bước kiểm tra giúp trả lời các câu hỏi: có bao nhiêu dòng? có những cột nào? kiểu dữ liệu là gì? có giá trị thiếu không?
\end{block}
\end{frame}

\begin{frame}[fragile]{Ví dụ kết quả: kiểm tra nhanh DataFrame}
\small
Với DataFrame sinh viên ở trên, một số lệnh kiểm tra cho kết quả như sau:
\begin{lstlisting}[style=py]
print(df.shape)
# (3, 3)

print(df.columns)
# Index(['Name', 'Age', 'GPA'], dtype='object')

print(df.dtypes)
# Name     object
# Age       int64
# GPA     float64
# dtype: object
\end{lstlisting}
\begin{block}{Cách đọc}
\cmd{shape = (3, 3)} nghĩa là bảng có 3 dòng và 3 cột. Kiểu \cmd{object} thường dùng cho dữ liệu chuỗi.
\end{block}
\end{frame}


\begin{frame}{Các lệnh chính: kiểm tra DataFrame}
\small
\begin{tabular}{p{0.28\linewidth}p{0.61\linewidth}}
\toprule
\textbf{Lệnh} & \textbf{Ý nghĩa} \\
\midrule
\cmd{df.head()} & Hiển thị các dòng đầu tiên. \\
\cmd{df.tail()} & Hiển thị các dòng cuối cùng. \\
\cmd{df.shape} & Trả về số dòng và số cột. \\
\cmd{df.columns} & Trả về nhãn cột. \\
\cmd{df.index} & Trả về nhãn dòng. \\
\cmd{df.dtypes} & Trả về kiểu dữ liệu của từng cột. \\
\cmd{df.info()} & Tóm tắt cấu trúc DataFrame. \\
\cmd{df.describe()} & Thống kê mô tả cho các cột số. \\
\bottomrule
\end{tabular}
\end{frame}

\begin{frame}[fragile]{Mini Exercise: tạo và kiểm tra dữ liệu}
\small
Tạo một DataFrame tên là \cmd{products} với các cột \cmd{Product}, \cmd{Price} và \cmd{Stock}.
\begin{lstlisting}[style=py]
products = pd.DataFrame({
    "Product": ["Laptop", "Mouse", "Keyboard"],
    "Price": [1200, 25, 45],
    "Stock": [5, 50, 30]
})

print(products.head())
print(products.shape)
print(products.dtypes)
\end{lstlisting}
\begin{alertblock}{Kiểm tra}
Tất cả các cột phải có cùng số lượng giá trị.
\end{alertblock}
\end{frame}

\begin{frame}[fragile]{Kết quả chạy code: DataFrame sản phẩm}
\small
Khi chạy phần mini exercise, sinh viên có thể kỳ vọng kết quả như sau:
\begin{lstlisting}[style=py]
    Product  Price  Stock
0    Laptop   1200      5
1     Mouse     25     50
2  Keyboard     45     30

print(products.shape)
# (3, 3)

print(products.dtypes)
# Product    object
# Price       int64
# Stock       int64
# dtype: object
\end{lstlisting}
\end{frame}


\begin{frame}{Quick Check: cơ bản về DataFrame}
\small
\quiztitle{Câu 1.} Lệnh nào tạo một DataFrame?
\begin{multicols}{2}
\begin{itemize}
  \item A. \cmd{pd.DataFrame()}
  \item B. \cmd{pd.SeriesFrame()}
  \item C. \cmd{np.DataFrame()}
  \item D. \cmd{df.create()}
\end{itemize}
\end{multicols}
\vspace{0.6em}
\quiztitle{Câu 2.} Lệnh nào trả về số dòng và số cột?
\begin{multicols}{2}
\begin{itemize}
  \item A. \cmd{df.shape}
  \item B. \cmd{df.mean()}
  \item C. \cmd{df.merge()}
  \item D. \cmd{df.plot()}
\end{itemize}
\end{multicols}
\end{frame}

\section{Series và lựa chọn dữ liệu}

\begin{frame}[fragile]{Tạo Series}
\small
\key{Series} là một mảng một chiều có nhãn.
\begin{lstlisting}[style=py]
s = pd.Series([10, 20, 30], index=["A", "B", "C"])
print(s)
print(s.index)
print(s.values)
print(s.dtype)
\end{lstlisting}
\begin{block}{Diễn giải}
Series giống như một cột có nhãn. Khi tính toán, Pandas căn chỉnh các Series theo nhãn index.
\end{block}
\end{frame}

\begin{frame}[fragile]{Kết quả chạy code: Series}
\small
\begin{lstlisting}[style=py]
A    10
B    20
C    30
dtype: int64

Index(['A', 'B', 'C'], dtype='object')
[10 20 30]
int64
\end{lstlisting}
\begin{block}{Cách đọc}
Cột bên trái là nhãn index; cột bên phải là giá trị của Series.
\end{block}
\end{frame}


\begin{frame}[fragile]{Căn chỉnh Series theo nhãn}
\small
\begin{lstlisting}[style=py]
x = pd.Series([10, 20, 30], index=["A", "B", "C"])
y = pd.Series([1, 2, 3], index=["B", "C", "D"])

result = x + y
print(result)
\end{lstlisting}
\begin{block}{Ý tưởng dự kiến}
Chỉ các nhãn khớp nhau mới được cộng trực tiếp. Các nhãn không khớp sẽ tạo ra giá trị thiếu.
\end{block}
\end{frame}

\begin{frame}[fragile]{Kết quả chạy code: căn chỉnh Series}
\small
\begin{lstlisting}[style=py]
A     NaN
B    21.0
C    32.0
D     NaN
dtype: float64
\end{lstlisting}
\begin{block}{Vì sao có NaN?}
Nhãn \cmd{B} và \cmd{C} xuất hiện trong cả hai Series nên được cộng. Nhãn \cmd{A} chỉ có trong \cmd{x}, còn \cmd{D} chỉ có trong \cmd{y}, nên kết quả là \cmd{NaN}.
\end{block}
\end{frame}


\begin{frame}[fragile]{Dữ liệu ví dụ cho lựa chọn, lọc và sắp xếp}
\small
Để các ví dụ \cmd{loc}, \cmd{iloc}, lọc và sắp xếp có kết quả rõ ràng, ta dùng bảng sinh viên sau.
\begin{lstlisting}[style=py]
df = pd.DataFrame({
    "StudentID": ["S01", "S02", "S03", "S04"],
    "Name": ["An", "Binh", "Chi", "Dung"],
    "Age": [20, 21, 19, 22],
    "GPA": [3.2, 3.6, 3.4, 3.8]
})

print(df)
\end{lstlisting}
\begin{lstlisting}[style=py]
  StudentID  Name  Age  GPA
0       S01    An   20  3.2
1       S02  Binh   21  3.6
2       S03   Chi   19  3.4
3       S04  Dung   22  3.8
\end{lstlisting}
\end{frame}

\begin{frame}[fragile]{Chọn cột}
\small
\begin{columns}[T]
\column{0.48\linewidth}
\textbf{Một cột trả về Series}
\begin{lstlisting}[style=py]
name = df["Name"]
print(name)
\end{lstlisting}
\column{0.48\linewidth}
\textbf{Nhiều cột trả về DataFrame}
\begin{lstlisting}[style=py]
subset = df[["Name", "GPA"]]
print(subset)
\end{lstlisting}
\end{columns}
\begin{block}{Ghi chú cú pháp}
Dùng một cặp ngoặc vuông để chọn một cột; dùng một list tên cột để chọn nhiều cột.
\end{block}
\end{frame}

\begin{frame}[fragile]{Kết quả chạy code: chọn cột}
\small
\begin{columns}[T]
\column{0.48\linewidth}
\textbf{Một cột: Series}
\begin{lstlisting}[style=py]
0      An
1    Binh
2     Chi
3    Dung
Name: Name, dtype: object
\end{lstlisting}
\column{0.48\linewidth}
\textbf{Nhiều cột: DataFrame}
\begin{lstlisting}[style=py]
   Name  GPA
0    An  3.2
1  Binh  3.6
2   Chi  3.4
3  Dung  3.8
\end{lstlisting}
\end{columns}
\end{frame}


\begin{frame}[fragile]{Chọn dòng bằng \cmd{loc} và \cmd{iloc}}
\small
\begin{columns}[T]
\column{0.48\linewidth}
\textbf{Chọn theo nhãn}
\begin{lstlisting}[style=py]
df = df.set_index("StudentID")
print(df.loc["S01"])
\end{lstlisting}
\column{0.48\linewidth}
\textbf{Chọn theo vị trí}
\begin{lstlisting}[style=py]
print(df.iloc[0])
print(df.iloc[0:2])
\end{lstlisting}
\end{columns}
\begin{alertblock}{Điểm khác biệt cốt lõi}
\cmd{loc} dùng nhãn; \cmd{iloc} dùng vị trí số nguyên.
\end{alertblock}
\end{frame}

\begin{frame}[fragile]{Kết quả chạy code: loc và iloc}
\small
Sau khi đặt \cmd{StudentID} làm index, \cmd{loc} có thể chọn theo mã sinh viên.
\begin{lstlisting}[style=py]
print(df.loc["S01"])
# Name     An
# Age      20
# GPA     3.2
# Name: S01, dtype: object

print(df.iloc[0:2])
#            Name  Age  GPA
# StudentID                
# S01          An   20  3.2
# S02        Binh   21  3.6
\end{lstlisting}
\end{frame}


\begin{frame}[fragile]{Chọn đồng thời dòng và cột}
\small
\begin{lstlisting}[style=py]
# chọn dòng và cột theo nhãn
selected = df.loc[["S01", "S02"], ["Name", "GPA"]]
print(selected)

# chọn dòng và cột theo vị trí
selected = df.iloc[0:2, [0, 2]]
print(selected)
\end{lstlisting}
\begin{block}{Cách đọc cú pháp}
Phần trước dấu phẩy chọn dòng; phần sau dấu phẩy chọn cột.
\end{block}
\end{frame}

\begin{frame}{Các lệnh chính: Series và lựa chọn dữ liệu}
\small
\begin{tabular}{p{0.30\linewidth}p{0.59\linewidth}}
\toprule
\textbf{Lệnh} & \textbf{Ý nghĩa} \\
\midrule
\cmd{pd.Series(data)} & Tạo mảng một chiều có nhãn. \\
\cmd{s.index} & Trả về nhãn của Series. \\
\cmd{s.values} & Trả về các giá trị bên trong. \\
\cmd{df["col"]} & Chọn một cột dưới dạng Series. \\
\cmd{df[["c1","c2"]]} & Chọn nhiều cột dưới dạng DataFrame. \\
\cmd{df.loc[label]} & Chọn theo nhãn dòng. \\
\cmd{df.iloc[position]} & Chọn theo vị trí số nguyên. \\
\bottomrule
\end{tabular}
\end{frame}

\begin{frame}{Quick Check: Series và lựa chọn}
\small
\quiztitle{Câu 1.} Pandas Series là:
\begin{itemize}
  \item A. một bảng hai chiều
  \item B. một mảng một chiều có nhãn
  \item C. một máy chủ cơ sở dữ liệu
  \item D. một thư viện vẽ biểu đồ
\end{itemize}
\vspace{0.5em}
\quiztitle{Câu 2.} Bộ chọn nào dựa trên nhãn?
\begin{multicols}{2}
\begin{itemize}
  \item A. \cmd{iloc}
  \item B. \cmd{loc}
  \item C. \cmd{shape}
  \item D. \cmd{head}
\end{itemize}
\end{multicols}
\end{frame}

\section{Lọc và sắp xếp dữ liệu}

\begin{frame}[fragile]{Lọc dữ liệu bằng điều kiện Boolean}
\small
Lọc các dòng thỏa mãn một điều kiện.
\begin{lstlisting}[style=py]
high_gpa = df[df["GPA"] >= 3.5]
print(high_gpa)
\end{lstlisting}
\begin{block}{Diễn giải}
\cmd{df["GPA"] >= 3.5} tạo một Series Boolean. Pandas giữ lại các dòng có giá trị \cmd{True}.
\end{block}
\end{frame}

\begin{frame}[fragile]{Kết quả chạy code: lọc dữ liệu}
\small
Với điều kiện \cmd{GPA >= 3.5}, Pandas chỉ giữ lại các dòng thỏa mãn điều kiện.
\begin{lstlisting}[style=py]
high_gpa = df[df["GPA"] >= 3.5]
print(high_gpa)

#            Name  Age  GPA
# StudentID                
# S02        Binh   21  3.6
# S04        Dung   22  3.8
\end{lstlisting}
\begin{block}{Cách đọc}
Hai sinh viên S02 và S04 có GPA lớn hơn hoặc bằng 3.5 nên xuất hiện trong kết quả.
\end{block}
\end{frame}


\begin{frame}[fragile]{Lọc với nhiều điều kiện}
\small
\begin{lstlisting}[style=py]
selected = df[
    (df["Age"] >= 20) &
    (df["GPA"] >= 3.5)
]
print(selected)
\end{lstlisting}
\begin{itemize}
  \item Dùng \cmd{&} cho AND.
  \item Dùng \cmd{|} cho OR.
  \item Dùng \cmd{~} cho NOT.
  \item Đặt từng điều kiện trong ngoặc tròn.
\end{itemize}
\end{frame}

\begin{frame}[fragile]{Sắp xếp DataFrame}
\small
\begin{lstlisting}[style=py]
# sắp xếp theo một cột
by_gpa = df.sort_values("GPA")

# sắp xếp giảm dần
by_gpa_desc = df.sort_values("GPA", ascending=False)

# sắp xếp theo nhiều cột
ordered = df.sort_values(
    ["Age", "GPA"],
    ascending=[True, False]
)
\end{lstlisting}
\begin{block}{Ví dụ kinh doanh}
Sắp xếp khách hàng theo doanh thu giảm dần, sau đó theo tần suất đặt hàng giảm dần.
\end{block}
\end{frame}

\begin{frame}[fragile]{Kết quả chạy code: lọc nhiều điều kiện và sắp xếp}
\small
\begin{lstlisting}[style=py]
selected = df[(df["Age"] >= 20) & (df["GPA"] >= 3.5)]
print(selected)

#            Name  Age  GPA
# StudentID                
# S02        Binh   21  3.6
# S04        Dung   22  3.8

print(df.sort_values("GPA", ascending=False))
#            Name  Age  GPA
# StudentID                
# S04        Dung   22  3.8
# S02        Binh   21  3.6
# S03         Chi   19  3.4
# S01          An   20  3.2
\end{lstlisting}
\end{frame}


\begin{frame}{Các lệnh chính: lọc và sắp xếp}
\small
\begin{tabular}{p{0.33\linewidth}p{0.56\linewidth}}
\toprule
\textbf{Lệnh} & \textbf{Ý nghĩa} \\
\midrule
\cmd{df[df["x"] > a]} & Giữ các dòng trong đó cột \cmd{x} lớn hơn \cmd{a}. \\
\cmd{(cond1) & (cond2)} & Kết hợp điều kiện bằng AND. \\
\cmd{(cond1) | (cond2)} & Kết hợp điều kiện bằng OR. \\
\cmd{~cond} & Phủ định một điều kiện. \\
\cmd{df.sort_values("x")} & Sắp xếp dòng theo một cột. \\
\cmd{ascending=False} & Sắp xếp giảm dần. \\
\bottomrule
\end{tabular}
\end{frame}

\begin{frame}{Quick Check: lọc và sắp xếp}
\small
\quiztitle{Câu 1.} Biểu thức nào lọc các dòng có \cmd{GPA >= 3.5}?
\begin{itemize}
  \item A. \cmd{df[df["GPA"] >= 3.5]}
  \item B. \cmd{df.sort_values("GPA")}
  \item C. \cmd{df.read_csv("GPA")}
  \item D. \cmd{df.merge("GPA")}
\end{itemize}
\vspace{0.5em}
\quiztitle{Câu 2.} Vì sao cần dùng ngoặc tròn khi kết hợp nhiều điều kiện lọc?
\end{frame}

\section{Nhập và xuất dữ liệu}

\begin{frame}[fragile]{Đọc và ghi tệp CSV}
\small
\begin{columns}[T]
\column{0.48\linewidth}
\textbf{Đọc CSV}
\begin{lstlisting}[style=py]
df = pd.read_csv("data.csv")

df = pd.read_csv(
    "data.csv",
    sep=",",
    header=0,
    encoding="utf-8"
)
\end{lstlisting}
\column{0.48\linewidth}
\textbf{Ghi CSV}
\begin{lstlisting}[style=py]
df.to_csv(
    "output.csv",
    index=False
)
\end{lstlisting}
\end{columns}
\begin{alertblock}{Quan trọng}
\cmd{index=False} giúp tránh xuất index của DataFrame thành một cột thừa.
\end{alertblock}
\end{frame}

\begin{frame}[fragile]{Đọc và ghi tệp Excel}
\small
\begin{columns}[T]
\column{0.48\linewidth}
\textbf{Đọc Excel}
\begin{lstlisting}[style=py]
df = pd.read_excel(
    "data.xlsx",
    sheet_name="Sheet1"
)
\end{lstlisting}
\column{0.48\linewidth}
\textbf{Ghi Excel}
\begin{lstlisting}[style=py]
df.to_excel(
    "output.xlsx",
    index=False
)
\end{lstlisting}
\end{columns}
\begin{block}{Ghi chú giảng dạy}
Excel rất phổ biến trong bối cảnh kinh doanh, nhưng sau khi nhập dữ liệu vẫn cần kiểm tra kiểu dữ liệu và giá trị thiếu.
\end{block}
\end{frame}

\begin{frame}[fragile]{JSON và tệp văn bản}
\small
\begin{columns}[T]
\column{0.48\linewidth}
\textbf{JSON}
\begin{lstlisting}[style=py]
df = pd.read_json("data.json")

df.to_json(
    "output.json",
    orient="records"
)
\end{lstlisting}
\column{0.48\linewidth}
\textbf{Tệp văn bản có phân tách}
\begin{lstlisting}[style=py]
df = pd.read_csv(
    "data.txt",
    sep="\t"
)
\end{lstlisting}
\end{columns}
\begin{block}{Ý tưởng}
Nhiều tệp văn bản có cấu trúc có thể được đọc bằng \cmd{read_csv()} bằng cách thay đổi ký tự phân tách.
\end{block}
\end{frame}

\begin{frame}{Các lệnh chính: nhập và xuất dữ liệu}
\small
\begin{tabular}{p{0.30\linewidth}p{0.59\linewidth}}
\toprule
\textbf{Lệnh} & \textbf{Vai trò} \\
\midrule
\cmd{pd.read_csv()} & Đọc tệp CSV hoặc tệp văn bản có phân tách. \\
\cmd{df.to_csv()} & Ghi dữ liệu ra tệp CSV. \\
\cmd{pd.read_excel()} & Đọc tệp Excel. \\
\cmd{df.to_excel()} & Ghi dữ liệu ra tệp Excel. \\
\cmd{pd.read_json()} & Đọc tệp JSON. \\
\cmd{df.to_json()} & Ghi dữ liệu ra tệp JSON. \\
\bottomrule
\end{tabular}
\end{frame}

\begin{frame}{Quick Check: nhập và xuất dữ liệu}
\small
\quiztitle{Câu 1.} Hàm nào dùng để đọc tệp CSV?
\begin{multicols}{2}
\begin{itemize}
  \item A. \cmd{pd.read_csv()}
  \item B. \cmd{pd.open_csv()}
  \item C. \cmd{pd.load_table_only()}
  \item D. \cmd{df.csv_read()}
\end{itemize}
\end{multicols}
\vspace{0.6em}
\quiztitle{Câu 2.} Tham số nào thường dùng để tránh xuất index ra tệp?
\begin{multicols}{2}
\begin{itemize}
  \item A. \cmd{index=False}
  \item B. \cmd{index=True}
  \item C. \cmd{header=None}
  \item D. \cmd{drop_index=True}
\end{itemize}
\end{multicols}
\end{frame}

\section{Làm sạch dữ liệu}

\begin{frame}{Vì sao làm sạch dữ liệu quan trọng?}
\small
Dữ liệu thực tế có thể chứa:
\begin{itemize}
  \item giá trị thiếu;
  \item dòng trùng lặp;
  \item kiểu dữ liệu không nhất quán;
  \item cột rỗng hoặc không liên quan;
  \item văn bản không thống nhất, thừa khoảng trắng hoặc khác chữ hoa/chữ thường;
  \item dữ liệu số bị lẫn với văn bản.
\end{itemize}
\begin{block}{Mục tiêu}
Làm sạch dữ liệu giúp tăng độ chính xác, tính nhất quán và khả năng sử dụng trước khi phân tích hoặc xây dựng mô hình.
\end{block}
\end{frame}

\begin{frame}[fragile]{Ví dụ dữ liệu bẩn và kết quả kiểm tra}
\small
\begin{lstlisting}[style=py]
df = pd.DataFrame({
    "Name": [" An ", "Binh", "Chi", "Chi"],
    "Age": [20, None, 22, 22],
    "Price": ["10", "20", "unknown", "30"]
})

print(df)
print(df.isna().sum())
\end{lstlisting}
\begin{lstlisting}[style=py]
   Name   Age    Price
0    An   20.0       10
1  Binh    NaN       20
2   Chi   22.0  unknown
3   Chi   22.0       30

Name     0
Age      1
Price    0
dtype: int64
\end{lstlisting}
\end{frame}


\begin{frame}[fragile]{Phát hiện giá trị thiếu}
\small
\begin{lstlisting}[style=py]
print(df.isna())
print(df.isna().sum())

missing_by_column = df.isna().sum()
total_missing = df.isna().sum().sum()
\end{lstlisting}
\begin{block}{Diễn giải}
\cmd{isna()} nhận diện các ô bị thiếu; \cmd{sum()} đếm số lượng giá trị thiếu theo cột hoặc toàn bộ DataFrame.
\end{block}
\end{frame}

\begin{frame}[fragile]{Xóa hoặc điền giá trị thiếu}
\small
\begin{columns}[T]
\column{0.48\linewidth}
\textbf{Xóa giá trị thiếu}
\begin{lstlisting}[style=py]
rows_complete = df.dropna()
cols_complete = df.dropna(axis=1)
\end{lstlisting}
\column{0.48\linewidth}
\textbf{Điền giá trị thiếu}
\begin{lstlisting}[style=py]
df["Age"] = df["Age"].fillna(
    df["Age"].mean()
)

df["City"] = df["City"].fillna("Unknown")
\end{lstlisting}
\end{columns}
\begin{alertblock}{Câu hỏi quyết định}
Nên xóa, suy diễn hay gắn nhãn riêng cho giá trị thiếu? Câu trả lời phụ thuộc vào mục tiêu phân tích.
\end{alertblock}
\end{frame}

\begin{frame}[fragile]{Dòng trùng lặp và kiểu dữ liệu}
\small
\begin{columns}[T]
\column{0.48\linewidth}
\textbf{Dòng trùng lặp}
\begin{lstlisting}[style=py]
duplicate_mask = df.duplicated()
duplicate_count = duplicate_mask.sum()
df = df.drop_duplicates()
\end{lstlisting}
\column{0.48\linewidth}
\textbf{Kiểu dữ liệu}
\begin{lstlisting}[style=py]
print(df.dtypes)

df["Quantity"] = df["Quantity"].astype(int)

df["Price"] = pd.to_numeric(
    df["Price"],
    errors="coerce"
)
\end{lstlisting}
\end{columns}
\end{frame}

\begin{frame}[fragile]{Làm sạch chuỗi văn bản}
\small
Các cột văn bản thường cần được chuẩn hóa.
\begin{lstlisting}[style=py]
df["Name"] = df["Name"].str.strip()
df["Name"] = df["Name"].str.lower()

df["City"] = df["City"].str.replace(
    "HN",
    "Hanoi"
)
\end{lstlisting}
\begin{block}{Vấn đề thường gặp}
\cmd{" An "}, \cmd{"AN"} và \cmd{"an"} có thể cùng chỉ một tên, nhưng là các chuỗi khác nhau nếu chưa được làm sạch.
\end{block}
\end{frame}

\begin{frame}[fragile]{Kết quả chạy code: làm sạch dữ liệu}
\small
\begin{lstlisting}[style=py]
df["Age"] = df["Age"].fillna(df["Age"].mean())
df["Name"] = df["Name"].str.strip().str.lower()
df["Price"] = pd.to_numeric(df["Price"], errors="coerce")
print(df)
\end{lstlisting}
\begin{lstlisting}[style=py]
   Name   Age  Price
0    an  20.0   10.0
1  binh  21.33  20.0
2   chi  22.0    NaN
3   chi  22.0   30.0
\end{lstlisting}
\begin{block}{Cách đọc}
\cmd{unknown} được chuyển thành \cmd{NaN}; tên được bỏ khoảng trắng và đưa về chữ thường.
\end{block}
\end{frame}


\begin{frame}{Các lệnh chính: làm sạch dữ liệu}
\small
\begin{tabular}{p{0.32\linewidth}p{0.57\linewidth}}
\toprule
\textbf{Lệnh} & \textbf{Ý nghĩa} \\
\midrule
\cmd{df.isna()} & Phát hiện giá trị thiếu. \\
\cmd{df.isna().sum()} & Đếm giá trị thiếu. \\
\cmd{df.dropna()} & Xóa quan sát bị thiếu dữ liệu. \\
\cmd{df.fillna(value)} & Thay thế giá trị thiếu. \\
\cmd{df.duplicated()} & Phát hiện dòng trùng lặp. \\
\cmd{df.drop_duplicates()} & Xóa dòng trùng lặp. \\
\cmd{df.astype(type)} & Chuyển kiểu dữ liệu. \\
\cmd{pd.to_numeric()} & Chuyển giá trị sang dạng số. \\
\cmd{Series.str.*} & Áp dụng các phương thức xử lý chuỗi. \\
\bottomrule
\end{tabular}
\end{frame}

\begin{frame}{Quick Check: làm sạch dữ liệu}
\small
\quiztitle{Câu 1.} Phương thức nào xóa các dòng có giá trị thiếu?
\begin{multicols}{2}
\begin{itemize}
  \item A. \cmd{dropna()}
  \item B. \cmd{fillna()}
  \item C. \cmd{duplicated()}
  \item D. \cmd{astype()}
\end{itemize}
\end{multicols}
\vspace{0.6em}
\quiztitle{Câu 2.} Phương thức nào xóa các dòng trùng lặp?
\begin{multicols}{2}
\begin{itemize}
  \item A. \cmd{drop_duplicates()}
  \item B. \cmd{dropna()}
  \item C. \cmd{sort_values()}
  \item D. \cmd{merge()}
\end{itemize}
\end{multicols}
\end{frame}

\section{Thao tác và tổng hợp dữ liệu}

\begin{frame}[fragile]{Dữ liệu ví dụ cho thao tác và tổng hợp}
\small
\begin{lstlisting}[style=py]
df = pd.DataFrame({
    "Region": ["North", "North", "South", "South", "East"],
    "Product": ["A", "B", "A", "B", "A"],
    "Quantity": [2, 3, 5, 4, 1],
    "Price": [10, 20, 8, 15, 12],
    "Sales": [20, 60, 40, 60, 12]
})
print(df)
\end{lstlisting}
\begin{lstlisting}[style=py]
  Region Product  Quantity  Price  Sales
0  North       A         2     10     20
1  North       B         3     20     60
2  South       A         5      8     40
3  South       B         4     15     60
4   East       A         1     12     12
\end{lstlisting}
\end{frame}

\begin{frame}[fragile]{Tạo cột tính toán}
\small
\begin{lstlisting}[style=py]
df["Revenue"] = df["Quantity"] * df["Price"]

total_revenue = df["Revenue"].sum()
print(total_revenue)
\end{lstlisting}
\begin{block}{Ví dụ kinh doanh}
Doanh thu, biên lợi nhuận, tỷ lệ giảm giá, giá trị vòng đời khách hàng và độ trễ giao hàng thường có thể được tạo từ các cột hiện có.
\end{block}
\end{frame}

\begin{frame}[fragile]{Kết quả chạy code: cột tính toán}
\small
\begin{lstlisting}[style=py]
df["Revenue"] = df["Quantity"] * df["Price"]
print(df[["Product", "Quantity", "Price", "Revenue"]])
print(df["Revenue"].sum())
\end{lstlisting}
\begin{lstlisting}[style=py]
  Product  Quantity  Price  Revenue
0       A         2     10       20
1       B         3     20       60
2       A         5      8       40
3       B         4     15       60
4       A         1     12       12

192
\end{lstlisting}
\end{frame}


\begin{frame}[fragile]{Áp dụng hàm}
\small
\begin{columns}[T]
\column{0.48\linewidth}
\textbf{Dùng \cmd{map()}}
\begin{lstlisting}[style=py]
df["Status"] = df["Score"].map(
    lambda x: "Pass" if x >= 5 else "Fail"
)
\end{lstlisting}
\column{0.48\linewidth}
\textbf{Dùng \cmd{apply()}}
\begin{lstlisting}[style=py]
df["Squared"] = df["Value"].apply(
    lambda x: x ** 2
)
\end{lstlisting}
\end{columns}
\begin{alertblock}{Thực hành tốt}
Ưu tiên các phép toán vector hóa trên cột khi có thể; dùng \cmd{apply()} khi cần logic tùy chỉnh theo dòng hoặc theo phần tử.
\end{alertblock}
\end{frame}

\begin{frame}[fragile]{Chuẩn hóa và chuẩn hóa Z-score}
\small
\begin{lstlisting}[style=py]
# Min-max normalization
df["Normalized"] = (
    (df["Value"] - df["Value"].min()) /
    (df["Value"].max() - df["Value"].min())
)

# Z-score standardization
df["Z"] = (
    (df["Value"] - df["Value"].mean()) /
    df["Value"].std()
)
\end{lstlisting}
\begin{block}{Diễn giải}
Min-max normalization đưa dữ liệu về một khoảng cố định; Z-score biểu diễn giá trị theo đơn vị độ lệch chuẩn.
\end{block}
\end{frame}

\begin{frame}[fragile]{Thống kê mô tả}
\small
\begin{lstlisting}[style=py]
print(df.describe())

print(df["Sales"].mean())
print(df["Sales"].median())
print(df["Sales"].min())
print(df["Sales"].max())
print(df["Sales"].std())
\end{lstlisting}
\begin{block}{Vai trò phân tích}
Thống kê mô tả giúp nhận biết thang đo, xu hướng trung tâm, độ biến thiên và các giá trị bất thường có thể tồn tại.
\end{block}
\end{frame}

\begin{frame}[fragile]{Gom nhóm với \cmd{groupby()}}
\small
\begin{lstlisting}[style=py]
region_total = df.groupby("Region")["Sales"].sum()
print(region_total)

region_stats = (
    df.groupby("Region")["Sales"]
      .agg(["count", "sum", "mean", "min", "max"])
)
print(region_stats)
\end{lstlisting}
\begin{block}{Split-apply-combine}
Pandas chia dữ liệu thành các nhóm, áp dụng phép tổng hợp cho từng nhóm, sau đó kết hợp kết quả.
\end{block}
\end{frame}

\begin{frame}[fragile]{Kết quả chạy code: groupby và agg}
\small
\begin{lstlisting}[style=py]
region_total = df.groupby("Region")["Sales"].sum()
print(region_total)
\end{lstlisting}
\begin{lstlisting}[style=py]
Region
East      12
North     80
South    100
Name: Sales, dtype: int64
\end{lstlisting}
\begin{lstlisting}[style=py]
region_stats = df.groupby("Region")["Sales"].agg(["count", "sum", "mean"])
print(region_stats)
\end{lstlisting}
\begin{lstlisting}[style=py]
        count  sum  mean
Region                  
East        1   12  12.0
North       2   80  40.0
South       2  100  50.0
\end{lstlisting}
\end{frame}


\begin{frame}[fragile]{Gom nhóm theo nhiều cột}
\small
\begin{lstlisting}[style=py]
summary = (
    df.groupby(["Region", "Product"])["Sales"]
      .sum()
)
print(summary)
\end{lstlisting}
\begin{block}{Diễn giải kinh doanh}
Cách này giúp trả lời các câu hỏi như: sản phẩm nào đóng góp doanh thu lớn nhất tại từng khu vực?
\end{block}
\end{frame}

\begin{frame}{Các lệnh chính: thao tác và tổng hợp}
\small
\begin{tabular}{p{0.32\linewidth}p{0.57\linewidth}}
\toprule
\textbf{Lệnh} & \textbf{Ý nghĩa} \\
\midrule
\cmd{df["new"] = ...} & Tạo hoặc biến đổi một cột. \\
\cmd{Series.map()} & Ánh xạ giá trị sang giá trị mới. \\
\cmd{Series.apply()} & Áp dụng hàm lên các phần tử của Series. \\
\cmd{df.describe()} & Thống kê mô tả. \\
\cmd{df.groupby()} & Gom nhóm quan sát. \\
\cmd{.agg()} & Áp dụng nhiều hàm tổng hợp. \\
\cmd{mean(), median(), std()} & Các thống kê số thường dùng. \\
\bottomrule
\end{tabular}
\end{frame}

\begin{frame}{Quick Check: thao tác dữ liệu}
\small
\quiztitle{Câu 1.} Phương thức nào dùng để gom nhóm quan sát theo category?
\begin{multicols}{2}
\begin{itemize}
  \item A. \cmd{groupby()}
  \item B. \cmd{dropna()}
  \item C. \cmd{astype()}
  \item D. \cmd{sort_index()}
\end{itemize}
\end{multicols}
\vspace{0.6em}
\quiztitle{Câu 2.} Hoàn thành công thức: \cmd{Revenue = Quantity ... Price}.
\end{frame}

\section{Kết hợp và định dạng lại dữ liệu}

\begin{frame}[fragile]{Nối DataFrame bằng concat}
\small
\begin{lstlisting}[style=py]
combined = pd.concat(
    [df1, df2],
    axis=0
)
\end{lstlisting}
\begin{itemize}
  \item \cmd{axis=0}: xếp chồng các dòng theo chiều dọc.
  \item \cmd{axis=1}: ghép các cột theo chiều ngang.
\end{itemize}
\begin{block}{Ví dụ}
Nối dữ liệu bán hàng tháng 1 và tháng 2 khi hai bảng có cùng cấu trúc cột.
\end{block}
\end{frame}

\begin{frame}[fragile]{Dữ liệu ví dụ cho merge}
\small
\begin{columns}[T]
\column{0.48\linewidth}
\begin{lstlisting}[style=py]
customers = pd.DataFrame({
    "CustomerID": [1, 2, 3],
    "Name": ["An", "Binh", "Chi"]
})
\end{lstlisting}
\column{0.48\linewidth}
\begin{lstlisting}[style=py]
orders = pd.DataFrame({
    "OrderID": [101, 102, 103],
    "CustomerID": [1, 1, 2],
    "Amount": [200, 150, 300]
})
\end{lstlisting}
\end{columns}
\begin{block}{Mục tiêu}
Ghép thông tin khách hàng với thông tin đơn hàng thông qua khóa chung \cmd{CustomerID}.
\end{block}
\end{frame}

\begin{frame}[fragile]{Ghép bảng bằng merge}
\small
\begin{lstlisting}[style=py]
result = pd.merge(
    customers,
    orders,
    on="CustomerID",
    how="inner"
)
\end{lstlisting}
\begin{itemize}
  \item \cmd{inner}: chỉ giữ các bản ghi khớp nhau.
  \item \cmd{left}: giữ toàn bộ bản ghi từ bảng bên trái.
  \item \cmd{right}: giữ toàn bộ bản ghi từ bảng bên phải.
  \item \cmd{outer}: giữ toàn bộ bản ghi từ cả hai bảng.
\end{itemize}
\end{frame}

\begin{frame}[fragile]{Kết quả chạy code: merge}
\small
\begin{lstlisting}[style=py]
result = pd.merge(customers, orders, on="CustomerID", how="inner")
print(result)
\end{lstlisting}
\begin{lstlisting}[style=py]
   CustomerID  Name  OrderID  Amount
0           1    An      101     200
1           1    An      102     150
2           2  Binh      103     300
\end{lstlisting}
\begin{block}{Cách đọc}
Khách hàng \cmd{Chi} không xuất hiện trong \cmd{inner merge} vì không có đơn hàng khớp trong bảng \cmd{orders}.
\end{block}
\end{frame}


\begin{frame}[fragile]{Pivot table}
\small
Pivot table tóm tắt các giá trị số theo các nhóm phân loại.
\begin{lstlisting}[style=py]
pivot = pd.pivot_table(
    sales,
    values="Revenue",
    index="Region",
    columns="Product",
    aggfunc="sum"
)
print(pivot)
\end{lstlisting}
\begin{block}{Diễn giải}
Các dòng là khu vực, các cột là sản phẩm, và mỗi ô chứa tổng doanh thu.
\end{block}
\end{frame}

\begin{frame}[fragile]{Kết quả chạy code: pivot table}
\small
Ví dụ nếu dữ liệu có doanh thu theo khu vực và sản phẩm, pivot table có thể có dạng:
\begin{lstlisting}[style=py]
Product     A     B
Region             
East     12.0   NaN
North    20.0  60.0
South    40.0  60.0
\end{lstlisting}
\begin{block}{Cách đọc}
Ô tại dòng \cmd{South}, cột \cmd{B} bằng 60, nghĩa là tổng doanh thu của sản phẩm B ở khu vực South là 60.
\end{block}
\end{frame}


\begin{frame}[fragile]{Định dạng lại dữ liệu: pivot và melt}
\small
\begin{columns}[T]
\column{0.48\linewidth}
\textbf{Từ long sang wide}
\begin{lstlisting}[style=py]
wide = df.pivot(
    index="Date",
    columns="Product",
    values="Sales"
)
\end{lstlisting}
\column{0.48\linewidth}
\textbf{Từ wide sang long}
\begin{lstlisting}[style=py]
long = pd.melt(
    wide.reset_index(),
    id_vars="Date"
)
\end{lstlisting}
\end{columns}
\begin{block}{Ý tưởng}
Dạng dữ liệu nên phù hợp với nhiệm vụ phân tích: trực quan hóa, tổng hợp, mô hình hóa hoặc báo cáo.
\end{block}
\end{frame}

\begin{frame}{Các lệnh chính: kết hợp và định dạng lại}
\small
\begin{tabular}{p{0.32\linewidth}p{0.57\linewidth}}
\toprule
\textbf{Lệnh} & \textbf{Ý nghĩa} \\
\midrule
\cmd{pd.concat()} & Kết hợp DataFrame theo một trục. \\
\cmd{pd.merge()} & Kết hợp DataFrame dựa trên cột khóa. \\
\cmd{how="inner"} & Chỉ giữ các bản ghi khớp. \\
\cmd{how="left"} & Giữ toàn bộ bản ghi từ bảng trái. \\
\cmd{df.pivot()} & Chuyển dữ liệu dạng long sang wide. \\
\cmd{pd.melt()} & Chuyển dữ liệu dạng wide sang long. \\
\cmd{pd.pivot_table()} & Tạo bảng tóm tắt. \\
\bottomrule
\end{tabular}
\end{frame}

\begin{frame}{Quick Check: kết hợp và định dạng lại}
\small
\quiztitle{Câu 1.} Hàm nào kết hợp các DataFrame bằng một cột khóa?
\begin{multicols}{2}
\begin{itemize}
  \item A. \cmd{pd.merge()}
  \item B. \cmd{pd.mean()}
  \item C. \cmd{pd.reshape()}
  \item D. \cmd{pd.filter_rows()}
\end{itemize}
\end{multicols}
\vspace{0.6em}
\quiztitle{Câu 2.} Pivot table dùng để làm gì?
\end{frame}

\section{Thao tác nâng cao}

\begin{frame}[fragile]{Tương quan}
\small
Tương quan đo mức độ liên hệ giữa các biến số.
\begin{lstlisting}[style=py]
print(df.corr(numeric_only=True))

print(
    df["Advertising"].corr(df["Sales"])
)
\end{lstlisting}
\begin{alertblock}{Quan trọng}
Tương quan không tự động chứng minh quan hệ nhân quả.
\end{alertblock}
\end{frame}

\begin{frame}[fragile]{Kết quả chạy code: tương quan và chuỗi thời gian}
\small
\begin{lstlisting}[style=py]
print(df[["Sales", "Quantity", "Price"]].corr())
\end{lstlisting}
\begin{lstlisting}[style=py]
             Sales  Quantity     Price
Sales     1.000000  0.620174  0.594089
Quantity  0.620174  1.000000 -0.272772
Price     0.594089 -0.272772  1.000000
\end{lstlisting}
\begin{block}{Cách đọc}
Ma trận tương quan đối xứng. Giá trị gần 1 biểu thị quan hệ tuyến tính dương mạnh hơn; tuy nhiên tương quan không chứng minh quan hệ nhân quả.
\end{block}
\end{frame}


\begin{frame}[fragile]{Trực quan hóa nhanh bằng Pandas}
\small
Các hàm vẽ của Pandas được xây dựng trên Matplotlib.
\begin{columns}[T]
\column{0.48\linewidth}
\begin{lstlisting}[style=py]
df.plot(x="Month", y="Sales", kind="line")

df.plot(x="Product", y="Sales", kind="bar")
\end{lstlisting}
\column{0.48\linewidth}
\begin{lstlisting}[style=py]
df["Sales"].plot(kind="hist")

df.plot(x="Advertising", y="Sales", kind="scatter")
\end{lstlisting}
\end{columns}
\end{frame}

\begin{frame}[fragile]{Quy trình xử lý chuỗi thời gian}
\small
\begin{lstlisting}[style=py]
df["Date"] = pd.to_datetime(df["Date"])
df = df.set_index("Date")
df = df.sort_index()

weekly = df["Sales"].resample("W").sum()
monthly = df["Sales"].resample("ME").sum()
\end{lstlisting}
\begin{block}{Ý tưởng chuỗi thời gian}
Sau khi chuyển cột ngày và đặt làm index, Pandas có thể tổng hợp dữ liệu theo các chu kỳ lịch.
\end{block}
\end{frame}

\begin{frame}[fragile]{Thành phần ngày tháng và trung bình trượt}
\small
\begin{columns}[T]
\column{0.48\linewidth}
\textbf{Thành phần ngày tháng}
\begin{lstlisting}[style=py]
df["Year"] = df.index.year
df["Month"] = df.index.month
df["Day"] = df.index.day
\end{lstlisting}
\column{0.48\linewidth}
\textbf{Trung bình trượt}
\begin{lstlisting}[style=py]
df["MovingAvg"] = (
    df["Sales"]
      .rolling(window=3)
      .mean()
)
\end{lstlisting}
\end{columns}
\end{frame}

\begin{frame}{Các lệnh chính: thao tác nâng cao}
\small
\begin{tabular}{p{0.32\linewidth}p{0.57\linewidth}}
\toprule
\textbf{Lệnh} & \textbf{Ý nghĩa} \\
\midrule
\cmd{df.corr()} & Tính ma trận tương quan. \\
\cmd{Series.corr()} & Tính tương quan giữa hai Series. \\
\cmd{df.plot()} & Trực quan hóa nhanh. \\
\cmd{pd.to_datetime()} & Chuyển dữ liệu sang kiểu datetime. \\
\cmd{df.set_index()} & Đặt một cột làm index. \\
\cmd{df.resample()} & Tổng hợp dữ liệu chuỗi thời gian theo khoảng thời gian. \\
\cmd{Series.rolling()} & Tính thống kê cửa sổ trượt. \\
\bottomrule
\end{tabular}
\end{frame}

\begin{frame}{Quick Check: thao tác nâng cao}
\small
\quiztitle{Câu 1.} Hàm nào chuyển một cột sang kiểu datetime?
\begin{multicols}{2}
\begin{itemize}
  \item A. \cmd{pd.to_datetime()}
  \item B. \cmd{pd.date_convert_only()}
  \item C. \cmd{df.datetime()}
  \item D. \cmd{pd.time_series()}
\end{itemize}
\end{multicols}
\vspace{0.5em}
\quiztitle{Câu 2.} Phương thức nào tính ma trận tương quan?
\begin{multicols}{2}
\begin{itemize}
  \item A. \cmd{corr()}
  \item B. \cmd{merge()}
  \item C. \cmd{dropna()}
  \item D. \cmd{pivot()}
\end{itemize}
\end{multicols}
\end{frame}

\section{Tóm tắt và thực hành}

\begin{frame}{Tóm tắt nội dung I}
\small
\begin{tabular}{p{0.27\linewidth}p{0.63\linewidth}}
\toprule
\textbf{Chủ đề} & \textbf{Ý chính} \\
\midrule
Pandas & Thư viện Python để xử lý và phân tích dữ liệu có cấu trúc. \\
Series & Mảng một chiều có nhãn. \\
DataFrame & Bảng dữ liệu hai chiều có nhãn. \\
Kiểm tra dữ liệu & Dùng \cmd{head()}, \cmd{shape}, \cmd{info()} và \cmd{describe()}. \\
Lựa chọn dữ liệu & Dùng ngoặc chọn cột, \cmd{loc} và \cmd{iloc}. \\
Lọc dữ liệu & Dùng điều kiện Boolean để giữ các dòng phù hợp. \\
\bottomrule
\end{tabular}
\end{frame}

\begin{frame}{Tóm tắt nội dung II}
\small
\begin{tabular}{p{0.27\linewidth}p{0.63\linewidth}}
\toprule
\textbf{Chủ đề} & \textbf{Ý chính} \\
\midrule
Nhập/xuất dữ liệu & Đọc và ghi CSV, Excel, JSON và tệp văn bản. \\
Làm sạch & Xử lý giá trị thiếu, trùng lặp, kiểu dữ liệu và chuỗi. \\
Thao tác & Tạo cột tính toán và biến đổi dữ liệu. \\
Tổng hợp & Dùng \cmd{groupby()} và \cmd{agg()} để tạo báo cáo tóm tắt. \\
Kết hợp & Dùng \cmd{concat()} và \cmd{merge()} cho nhiều bảng dữ liệu. \\
Nâng cao & Tương quan, trực quan hóa và thao tác chuỗi thời gian. \\
\bottomrule
\end{tabular}
\end{frame}

\begin{frame}{Thực hành: dữ liệu sinh viên}
\small
\begin{block}{Bài tập 1}
Tạo một DataFrame có ít nhất 10 sinh viên với các cột:
\begin{itemize}
  \item \cmd{StudentID}, \cmd{Name}, \cmd{Age}, \cmd{Major}, \cmd{GPA}.
\end{itemize}
Thực hiện:
\begin{enumerate}
  \item kiểm tra bằng \cmd{head()}, \cmd{shape}, \cmd{info()} và \cmd{describe()};
  \item lọc các sinh viên có \cmd{GPA >= 3.5};
  \item sắp xếp theo GPA giảm dần;
  \item chọn dòng bằng \cmd{loc} và \cmd{iloc}.
\end{enumerate}
\end{block}
\end{frame}

\begin{frame}{Thực hành: dữ liệu thiếu và làm sạch}
\small
\begin{block}{Bài tập 2}
Thêm giá trị thiếu vào tập dữ liệu sinh viên. Sau đó:
\begin{enumerate}
  \item đếm số giá trị thiếu;
  \item điền Age thiếu bằng Age trung bình;
  \item điền GPA thiếu bằng GPA trung vị;
  \item xóa các dòng thiếu Name.
\end{enumerate}
\end{block}
\begin{block}{Bài tập 3}
Tạo một dữ liệu lộn xộn có khoảng trắng trong tên, dòng trùng lặp và giá trị số bị lẫn văn bản. Làm sạch bằng \cmd{str.strip()}, \cmd{drop_duplicates()} và \cmd{pd.to_numeric()}.
\end{block}
\end{frame}

\begin{frame}{Thực hành: phân tích bán hàng}
\small
\begin{block}{Bài tập 4}
Tạo một tập dữ liệu bán hàng gồm \cmd{Date}, \cmd{Region}, \cmd{Product}, \cmd{Quantity} và \cmd{Price}.
Tạo:
\[
\text{Revenue} = \text{Quantity} \times \text{Price}.
\]
Sau đó tính:
\begin{enumerate}
  \item tổng doanh thu;
  \item doanh thu theo Region;
  \item doanh thu theo Product;
  \item doanh thu trung bình theo Region;
  \item pivot table Region $\times$ Product.
\end{enumerate}
\end{block}
\end{frame}

\begin{frame}{Thực hành: merge và chuỗi thời gian}
\small
\begin{block}{Bài tập 5: Merge}
Tạo \cmd{customers(CustomerID, Name, City)} và \cmd{orders(OrderID, CustomerID, Amount)}. Thực hiện inner merge, left merge và xác định khách hàng chưa có đơn hàng.
\end{block}
\begin{block}{Bài tập 6: Chuỗi thời gian}
Tạo dữ liệu bán hàng ngày trong ít nhất 30 ngày. Chuyển \cmd{Date} sang datetime, đặt làm index, tính tổng doanh thu theo tuần, tính trung bình trượt 7 ngày và vẽ line plot.
\end{block}
\end{frame}

\begin{frame}{Câu hỏi phản tư cuối bài}
\small
\begin{itemize}
  \item Series và DataFrame khác nhau như thế nào?
  \item Khi nào nên dùng \cmd{loc} và khi nào nên dùng \cmd{iloc}?
  \item Vì sao cần kiểm tra \cmd{shape}, \cmd{dtypes} và giá trị thiếu từ sớm?
  \item \cmd{concat()} khác \cmd{merge()} như thế nào?
  \item Vì sao \cmd{groupby()} quan trọng trong phân tích kinh doanh?
  \item Vì sao tương quan không tự động chứng minh quan hệ nhân quả?
\end{itemize}
\end{frame}

\begin{frame}{Đáp án gợi ý cho Quick Check}
\small
\begin{tabular}{p{0.42\linewidth}p{0.45\linewidth}}
\toprule
\textbf{Phần} & \textbf{Đáp án} \\
\midrule
Pandas là gì? & B; C \\
Cơ bản về DataFrame & A; A \\
Series và lựa chọn & B; B \\
Nhập và xuất dữ liệu & A; A \\
Làm sạch dữ liệu & A; A \\
Thao tác dữ liệu & A; Revenue = Quantity $\times$ Price \\
Kết hợp và định dạng lại & A; tóm tắt giá trị theo category \\
Thao tác nâng cao & A; A \\
\bottomrule
\end{tabular}
\end{frame}

\end{document}
